\documentclass[border=10pt]{standalone}
\usepackage[no-math]{fontspec}
\usepackage[brazil]{babel}
\usepackage{amsmath,amssymb}
\IfFileExists{C:/Windows/Fonts/Lato_Regular.ttf}{%
  \setmainfont{Lato}[
    Path = C:/Windows/Fonts/,
    UprightFont = Lato_Regular.ttf,
    BoldFont = Lato_Bold.ttf,
    ItalicFont = Lato_Italic.ttf,
    BoldItalicFont = Lato_BoldItalic.ttf
  ]
  \setsansfont{Lato}[
    Path = C:/Windows/Fonts/,
    UprightFont = Lato_Regular.ttf,
    BoldFont = Lato_Bold.ttf,
    ItalicFont = Lato_Italic.ttf,
    BoldItalicFont = Lato_BoldItalic.ttf
  ]
}{%
  \RequirePackage[default]{lato}
}
\renewcommand{\familydefault}{\sfdefault}

\usepackage{xcolor}
\usepackage{tikz}
\usepackage{pgfplots}
\pgfplotsset{compat=1.18}
\usetikzlibrary{arrows.meta,positioning,shapes.geometric,fit,backgrounds,calc,decorations.pathreplacing,patterns}

% Paleta institucional UEFS / Bioengenharia de Solos
\definecolor{bsNavy}{HTML}{1F4268}
\definecolor{bsNavyDark}{HTML}{152E4A}
\definecolor{bsGold}{HTML}{B38A3D}
\definecolor{bsGoldDark}{HTML}{8C6927}
\definecolor{bsTeal}{HTML}{0E7C7B}
\definecolor{bsForest}{HTML}{2D6A4F}
\definecolor{bsAlert}{HTML}{9E2A2B}
\definecolor{bsEarth}{HTML}{8C5A3C}
\definecolor{bsWater}{HTML}{2E86AB}
\definecolor{bsInk}{HTML}{26313D}
\definecolor{bsMuted}{HTML}{5B6573}

\definecolor{bgNavy}{HTML}{EDF2F7}
\definecolor{bgGold}{HTML}{F8F3E9}
\definecolor{bgTeal}{HTML}{EBF6F5}
\definecolor{bgForest}{HTML}{EEF5F0}
\definecolor{bgAlert}{HTML}{FDF2F2}

\tikzset{
  >=Stealth,
  bsBox/.style={
    draw=#1, fill=white, rounded corners=5pt, line width=1.15pt,
    inner sep=7pt, align=center, font=\footnotesize\sffamily, text=bsInk
  },
  bsTint/.style={
    draw=#1, fill=#1!8!white, rounded corners=5pt, line width=1.15pt,
    inner sep=7pt, align=center, font=\footnotesize\sffamily, text=bsInk
  },
  bsBadge/.style={
    fill=#1, text=white, rounded corners=9pt, inner xsep=8pt, inner ysep=3pt,
    font=\scriptsize\bfseries\sffamily
  },
  bsTitle/.style={font=\normalsize\bfseries\sffamily, text=bsNavy, align=center},
  bsSub/.style={font=\footnotesize\sffamily, text=bsMuted, align=center},
  bsArrow/.style={->, >=Stealth, line width=1.35pt},
  bsDArrow/.style={<->, >=Stealth, line width=1.35pt}
}


\begin{document}
\begin{tikzpicture}
  \hyphenpenalty=10000\exhyphenpenalty=10000

  % Título (bem acima dos cards superiores)
  \node[bsTitle] at (7.15, 8.60)
    {Zoneamento Integrado de Intervenções do Projeto Final na Vertente};
  \node[bsSub] at (7.15, 8.10)
    {Articulação topográfica da crista ao canal fluvial para controle contínuo de água e sedimentos};

  % ============================================================
  % PERFIL DA VERTENTE E CURSO D'ÁGUA
  % ============================================================
  \fill[bsEarth!14!white]
    (0.0, 0.0) -- (0.0, 4.95) -- (2.2, 4.95) -- (3.0, 4.55)
    .. controls (4.6, 4.0) and (6.2, 3.25) .. (7.8, 2.55)
    .. controls (9.2, 1.95) and (10.3, 1.45) .. (11.5, 1.15)
    -- (12.2, 0.55) -- (13.3, 0.55) -- (14.2, 0.95)
    -- (14.2, 0.0) -- cycle;

  % Água no canal fluvial (direita)
  \fill[bsWater!28!white]
    (11.85, 0.90) -- (14.05, 0.90) -- (13.3, 0.55) -- (12.2, 0.55) -- cycle;
  \draw[bsWater, line width=1.1pt] (11.85, 0.90) -- (14.05, 0.90);
  \node[font=\tiny\bfseries, text=bsWater] at (12.95, 0.72) {Canal Fluvial};

  % Linha do terreno
  \draw[bsEarth, line width=1.7pt]
    (0.0, 4.95) -- (2.2, 4.95) -- (3.0, 4.55)
    .. controls (4.6, 4.0) and (6.2, 3.25) .. (7.8, 2.55)
    .. controls (9.2, 1.95) and (10.3, 1.45) .. (11.5, 1.15)
    -- (12.2, 0.55) -- (13.3, 0.55) -- (14.2, 0.95);

  % Faixas verticais sutis de zoneamento no terreno
  \draw[bsMuted!40!white, dashed, line width=0.7pt] (3.4, 0.2) -- (3.4, 5.1);
  \draw[bsMuted!40!white, dashed, line width=0.7pt] (7.2, 0.2) -- (7.2, 5.1);
  \draw[bsMuted!40!white, dashed, line width=0.7pt] (10.6, 0.2) -- (10.6, 5.1);

  % ============================================================
  % SÍMBOLOS CONSTRUTIVOS SOBRE O PERFIL
  % ============================================================
  % Zona 1 (Crista): Canaleta Verde / Terraço
  \fill[bsWater!25!white] (1.6, 4.95) .. controls (1.9, 4.72) and (2.1, 4.72) .. (2.4, 4.95) -- cycle;
  \draw[bsWater, line width=1.3pt] (1.55, 4.95) .. controls (1.9, 4.70) and (2.1, 4.70) .. (2.45, 4.95);
  \draw[bsForest, line width=1.2pt] (1.7, 4.95) -- (1.65, 5.18);
  \draw[bsForest, line width=1.2pt] (2.3, 4.95) -- (2.35, 5.18);

  % Zona 2 (Meia-encosta): Biomanta + Fascinas + Dreno Verde
  \draw[bsForest, line width=2.6pt, opacity=0.7] (3.8, 4.25) -- (6.7, 3.05);
  \fill[bsForest, draw=bsForest!70!black] (4.4, 4.02) circle (0.11);
  \fill[bsForest, draw=bsForest!70!black] (5.5, 3.55) circle (0.11);
  \draw[bsWater, dashed, line width=1pt] (4.8, 3.75) -- (6.9, 2.72);

  % Zona 3 (Talvegue / Ravina): Paliçadas em série + bacia
  \draw[bsGoldDark, line width=2.2pt] (7.75, 2.58) -- (7.75, 2.92);
  \draw[bsGoldDark, line width=2.2pt] (8.85, 2.10) -- (8.85, 2.45);
  \draw[bsGoldDark, line width=2.2pt] (9.95, 1.65) -- (9.95, 2.00);
  \fill[bsGold!40!white] (7.75, 2.58) -- (8.4, 2.30) -- (7.75, 2.78) -- cycle;
  \fill[bsGold!40!white] (8.85, 2.10) -- (9.5, 1.82) -- (8.85, 2.30) -- cycle;

  % Zona 4 (Base / Pé & Margem Fluvial): Parede Krainer / Gabião / Enrocamento
  \fill[bsEarth!55!white, draw=bsEarth, line width=1pt] (11.15, 1.15) rectangle (11.85, 1.60);
  \draw[bsForest, line width=1.2pt] (11.35, 1.60) -- (11.20, 1.92);
  \draw[bsForest, line width=1.2pt] (11.65, 1.60) -- (11.80, 1.95);
  \fill[bsMuted!60!white, draw=bsInk] (11.95, 0.88) circle (0.13);
  \fill[bsMuted!55!white, draw=bsInk] (12.15, 0.72) circle (0.12);

  % Seta de fluxo controlado ao longo da vertente
  \draw[->, >=Stealth, bsWater, dashed, line width=1.1pt] (0.5, 5.22) -- (1.45, 5.08);

  % ============================================================
  % CARDS DE ZONEAMENTO (COM FOLGA DO SUBTÍTULO E ENTRE SI)
  % ============================================================
  \node[bsTint=bsWater, text width=3.15cm, minimum height=1.65cm] (z1) at (1.55, 6.45) {
    \vspace{0.18cm}
    \textbf{\color{bsWater}\scriptsize Zona 1 $\cdot$ Crista / Topo}\\[3pt]
    \scriptsize Terraços \& Canaleta Verde (desvio do escoamento)
  };
  \node[bsBadge=bsWater, font=\tiny\bfseries] at (z1.north) {CAPTAÇÃO MONTANTE};
  \draw[bsWater, line width=1pt, ->] (z1.south) -- (1.95, 5.15);

  \node[bsTint=bsForest, text width=3.35cm, minimum height=1.65cm] (z2) at (5.35, 6.45) {
    \vspace{0.18cm}
    \textbf{\color{bsForest}\scriptsize Zona 2 $\cdot$ Meia-Encosta}\\[3pt]
    \scriptsize Biomantas, Hidrossemeadura, Fascinas \& Dreno Verde
  };
  \node[bsBadge=bsForest, font=\tiny\bfseries] at (z2.north) {COBERTURA \& DRENAGEM};
  \draw[bsForest, line width=1pt, ->] (z2.south) -- (5.35, 3.85);

  \node[bsTint=bsGold, text width=3.25cm, minimum height=1.65cm] (z3) at (9.10, 5.30) {
    \vspace{0.18cm}
    \textbf{\color{bsGoldDark}\scriptsize Zona 3 $\cdot$ Talvegue / Ravina}\\[3pt]
    \scriptsize Paliçadas Permeáveis em Série \& Bacias de Captação
  };
  \node[bsBadge=bsGold, font=\tiny\bfseries] at (z3.north) {DISSIPAÇÃO EM DEGRAUS};
  \draw[bsGoldDark, line width=1pt, ->] (z3.south) -- (9.00, 2.60);

  \node[bsTint=bsEarth, text width=3.30cm, minimum height=1.65cm] (z4) at (12.55, 3.85) {
    \vspace{0.18cm}
    \textbf{\color{bsEarth}\scriptsize Zona 4 $\cdot$ Base \& Margem}\\[3pt]
    \scriptsize Parede Krainer, Gabião Vivo ou Enrocamento Vegetado
  };
  \node[bsBadge=bsEarth, font=\tiny\bfseries] at (z4.north) {CONTENÇÃO DE PÉ};
  \draw[bsEarth, line width=1pt, ->] (z4.south) -- (11.70, 2.05);

\end{tikzpicture}
\end{document}
